\documentclass[aspectratio=169,11pt]{beamer}
\usepackage[utf8]{inputenc}
\usepackage[T1]{fontenc}
\usepackage{lmodern}
\usepackage[spanish]{babel}
\usepackage{xurl}
\usetheme{Madrid}
\definecolor{itescaRed}{HTML}{8F1117}
\usecolortheme[named=itescaRed]{structure}
\setbeamertemplate{navigation symbols}{}
\title[Objetivos de vTaxi]{Objetivos de investigación de vTaxi}
\subtitle{Actividad 10 - Seminario I}
\author[M. J. de la Cruz Muñoz]{Martín Jonathan de la Cruz Muñoz\\\small Matrícula: 26130503}
\institute[ITESCA]{Instituto Tecnológico Superior de Cajeme\\Maestría en Gestión Administrativa\\Docente: Dra. Carla Olimpya Zapuche Moreno}
\titlegraphic{\includegraphics[height=0.65cm]{ITESCA/assets-itesca/web/logo-itesca-contorno.png}}
\hypersetup{pdfsubject={Actividad 10 - Seminario I}}
\date[04/10/2026]{4 de octubre de 2026}
\begin{document}
\begin{frame}\titlepage\end{frame}
\begin{frame}{Problema y pregunta general}\small Los borradores regulatorios de vTaxi no acreditan una ruta validada de requisitos, responsables y evidencias para su formalización como plataforma de taxis concesionados en Nuevo León.\par\medskip ¿Cómo se relaciona el nivel de gestión del cumplimiento normativo con el grado de preparación administrativa de vTaxi para su formalización como plataforma de taxis concesionados en Nuevo León durante septiembre a noviembre de 2026?\end{frame}
\begin{frame}{Objetivo general}\large Analizar la relación entre el nivel de gestión del cumplimiento normativo y el grado de preparación administrativa de vTaxi para su formalización como plataforma de taxis concesionados en Nuevo León durante septiembre a noviembre de 2026, mediante el contraste documental de requisitos, responsabilidades y evidencias del proyecto.\end{frame}
\begin{frame}{Objetivo específico 1}\small\textbf{Pregunta}\par ¿Qué requisitos y competencias institucionales corresponden al modelo de taxis concesionados de vTaxi y cuáles pertenecen a la modalidad de transporte privado mediante empresas de redes de transporte?\par\medskip\textbf{Objetivo}\par Identificar y clasificar los requisitos y las competencias institucionales aplicables al modelo de taxis concesionados de vTaxi, distinguiéndolos de los correspondientes al transporte privado mediante empresas de redes de transporte.\par\medskip\textbf{Procedimiento}\par Revisar las disposiciones aplicables y contrastarlas con el modelo descrito en los borradores de vTaxi. Registrar modalidad, competencia y localizador de fuente; separar expresamente taxis concesionados y transporte privado mediante empresas de redes de transporte.\par\medskip\textbf{Evidencia}\par Catálogo de requisitos con modalidad, autoridad, fuente y evidencia exigible. Permitirá distinguir cada requisito respaldado por una disposición aplicable de las recomendaciones internas y de las cuestiones que requieren aclaración institucional.\end{frame}
\begin{frame}{Objetivo específico 2}\small\textbf{Pregunta}\par ¿Qué diferencias existen entre los requisitos aplicables y las evidencias jurídicas, técnicas y operativas disponibles en los borradores del proyecto?\par\medskip\textbf{Objetivo}\par Diagnosticar las brechas entre los requisitos aplicables y las evidencias jurídicas, técnicas y operativas disponibles en los borradores del proyecto, según su integridad, coherencia y respaldo documental.\par\medskip\textbf{Procedimiento}\par Cotejar cada requisito clasificado con la documentación interna disponible. Registrar suficiencia, coherencia e información no localizada, conservando la referencia al documento y evitando presentar ejemplos de los borradores como datos comprobados del proyecto.\par\medskip\textbf{Evidencia}\par Diagnóstico de brechas por requisito, con documento localizado, inconsistencia y evidencia faltante. Cada hallazgo deberá remitir a un documento identificable; la ausencia de información no se registrará como cumplimiento ni incumplimiento acreditado.\end{frame}
\begin{frame}{Objetivo específico 3}\small\textbf{Pregunta}\par ¿Cómo se vinculan la asignación de responsables, la verificación de requisitos y el seguimiento documental con la integridad del expediente y la resolución de observaciones internas?\par\medskip\textbf{Objetivo}\par Examinar la correspondencia entre la asignación de responsables, la verificación de requisitos y el seguimiento documental, por una parte, y la integridad del expediente y la resolución de observaciones internas, por otra, mediante una revisión trazable de los componentes del caso.\par\medskip\textbf{Procedimiento}\par Comparar, dentro de cada componente del expediente, los controles de asignación y seguimiento con su respaldo documental y las observaciones registradas. Explicar coincidencias y diferencias como correspondencias del caso, no como correlaciones estadísticas poblacionales.\par\medskip\textbf{Evidencia}\par Matriz de contraste entre controles de gestión y estado del expediente. Para cada componente se registrarán responsable documentado, seguimiento, evidencia disponible y estado de observaciones, sin inferir efectos causales a partir del caso.\end{frame}
\begin{frame}{Objetivo específico 4}\small\textbf{Pregunta}\par ¿Qué acciones de gestión deben priorizarse para fortalecer la preparación del expediente de vTaxi sin confundir capacitación de conductores con autorización de la plataforma?\par\medskip\textbf{Objetivo}\par Proponer una ruta priorizada de gestión para fortalecer la preparación administrativa del expediente de vTaxi, con acciones, responsables propuestos y evidencias de atención, diferenciando la capacitación de conductores de la autorización de la plataforma.\par\medskip\textbf{Procedimiento}\par Ordenar las brechas identificadas por fundamento y dependencia administrativa. Proponer acciones y responsables funcionales, señalando la evidencia necesaria para cerrar cada observación y distinguiendo formación de conductores, concesiones y autorización de la plataforma.\par\medskip\textbf{Evidencia}\par Ruta priorizada de atención con acción, responsable propuesto, dependencia y criterio documental de cierre. Las prioridades se justificarán por pertinencia normativa y secuencia de integración; no representarán compromisos aceptados ni autorización de la plataforma.\end{frame}
\begin{frame}{Círculo de Covey}\small\begin{block}{Importante / No urgente}\textbf{Problema general}\par Los borradores regulatorios de vTaxi no acreditan una ruta validada de requisitos, responsables y evidencias para su formalización como plataforma de taxis concesionados en Nuevo León.\par\medskip\textbf{Objetivo general}\par Analizar la relación entre el nivel de gestión del cumplimiento normativo y el grado de preparación administrativa de vTaxi para su formalización como plataforma de taxis concesionados en Nuevo León durante septiembre a noviembre de 2026, mediante el contraste documental de requisitos, responsabilidades y evidencias del proyecto.\end{block}\end{frame}
\begin{frame}{Coherencia metodológica}\small La secuencia parte de identificar el marco aplicable, diagnostica las brechas, examina la correspondencia entre controles y evidencias y culmina en una propuesta priorizada. Así, cada objetivo específico responde a una pregunta ya formulada y aporta un resultado necesario para el objetivo general. La medibilidad se entiende como posibilidad de comprobar productos y hallazgos con documentos identificables; no requiere inventar porcentajes de mejora ni tamaños de muestra. Esta interpretación mantiene la exigencia de acciones observables y coherencia metodológica del material de objetivos (Instituto Tecnológico Superior de Cajeme, s. f.).\end{frame}
\begin{frame}{Conclusiones 1}\small La congruencia de los objetivos depende de que cada meta específica aporte evidencia a la pregunta general. En vTaxi, clasificar requisitos permite establecer el referente del diagnóstico; identificar brechas muestra qué está documentado; examinar controles y evidencias permite interpretar la preparación; y priorizar acciones transforma esa interpretación en una propuesta administrativa. La secuencia satisface la orientación hacia metas observables del material metodológico (Instituto Tecnológico Superior de Cajeme, s. f.).\end{frame}
\begin{frame}{Conclusiones 2}\small Considero que la aportación más útil del diseño es separar la preparación verificable de la expectativa de autorización: la primera puede estudiarse mediante documentos, mientras la segunda corresponde a la autoridad competente. En la gestión profesional, esa distinción evita tratar un borrador como prueba de cumplimiento o confundir capacitación con permiso. Por ello, la propuesta deberá conservar la incertidumbre donde falte sustento y justificar cada prioridad. Los objetivos delimitan el trabajo por realizar; no acreditan que el diagnóstico, la formalización o sus beneficios ya se hayan conseguido.\end{frame}
\begin{frame}{Fuente consultada}\small H. Congreso del Estado de Nuevo León. (2026). Ley de Movilidad Sostenible, de Accesibilidad y Seguridad Vial para el Estado de Nuevo León (última reforma publicada el 22 de mayo de 2026; arts. 82, 83, 99 y 100). \url{https://www.hcnl.gob.mx/trabajo_legislativo/leyes/leyes/ley_de_movilidad_sostenible_y_accesibilidad_para_el_estado_de_nuevo_leon/}\end{frame}
\begin{frame}{Fuente consultada}\small Instituto de Capacitación y Educación para el Trabajo del Estado de Nuevo León. (s. f.). ICET. Recuperado el 27 de septiembre de 2026, de \url{https://icetnl.mx/}\end{frame}
\begin{frame}{Fuente consultada}\small Instituto Tecnológico Superior de Cajeme. (s. f.). Formulación de objetivos general y específicos [Material didáctico de Seminario I, Unidad 3, apartado 3.3]. \url{https://cursos3.e-itesca.edu.mx/mod/resource/view.php?id=2897}\end{frame}
\begin{frame}{Asistencia de inteligencia artificial}\small Se utilizó GitHub Copilot para revisar la coherencia de los objetivos previamente formulados, redactar apoyos metodológicos y preparar los documentos. La asistencia automatizada no acredita revisión o aprobación docente ni sustituye la valoración crítica y la revisión final del estudiante.\end{frame}
\end{document}
